% TOP_PROBLEM: {"id": 2689, "title": "Kirby Problem 1.30", "category_id": 7, "category": {"id": 7, "name": "topology", "display_name": "Topology", "description": "Properties preserved under continuous deformations.", "slug": "topology", "order_index": 7, "created_at": "2026-07-31T15:26:25.671Z"}, "status": "open", "problem_number": "KP-1.30", "set_id": 11}
% FIRST_POSTED: 2026-10-01
% ENTRY_KIND: statement_only
% UnsolvedMath ID 2689; source catalogue code KP-1.30
% !TeX program = lualatex
% Definition and sources only; this file is not an LLM solution attempt.
\documentclass[11pt,a4paper]{article}
\usepackage[margin=25mm]{geometry}
\usepackage{fontspec}
\setmainfont{lmroman10-regular.otf}[BoldFont=lmroman10-bold.otf,
 ItalicFont=lmroman10-italic.otf,BoldItalicFont=lmroman10-bolditalic.otf]
\usepackage{amsmath,amssymb,mathtools}
\usepackage{unicode-math}
\setmathfont{latinmodern-math.otf}
\AtBeginDocument{\let\setminus\smallsetminus}
\usepackage{xurl,hyperref}
\hypersetup{colorlinks=true,urlcolor=blue}
\setcounter{secnumdepth}{0}
\setlength{\parindent}{0pt}
\setlength{\parskip}{5pt}
\setlength{\emergencystretch}{3em}
\begin{document}
\section*{Kirby Problem 1.30}\label{2689}
{\small UnsolvedMath ID 2689 · KP-1.30 · Topology · Kirby's Problems in Low-Dimensional Topology}
\subsection{Definitions and mathematical statement}
Let $K\subset S^3$ be a smooth knot. Use unreduced integral Khovanov homology
$\mathrm{Kh}^{i,j}(K;\mathbb Z)$, a bigraded finitely generated abelian group.
It is obtained from the cube of the two smoothings at each crossing of a
diagram: a resolution with $r$ circles receives the tensor power
$A^{\otimes r}$, where $A=\mathbb Z[X]/(X^2)$.
Edges of the cube use multiplication for merging circles and comultiplication
$\Delta(1)=1\otimes X+X\otimes1$, $\Delta(X)=X\otimes X$ for splitting them.
The usual signs and grading shifts give the link-invariant homology.

The conjecture states that, whenever $K$ is not isotopic to the unknot, there
exist gradings $i,j$ and a nonzero element $x$ with
\[
x\in\mathrm{Kh}^{i,j}(K;\mathbb Z),\qquad 2x=0.
\]
Thus the integral group must contain an element of order exactly two. It is
essential to retain integer coefficients: working over a field alone loses the
torsion being asked for. The assertion applies to all nontrivial knots, without
an alternating, positive, or crossing-number restriction.
\subsection{Short English statement}
Does every nontrivial knot have an element of order two in its integral Khovanov homology?
\subsection{Sources}
\begin{itemize}
\item[S1] ULAM.AI and the UnsolvedMath Contributors, supplied problems.json, record 2689 (KP-1.30), Kirby's Problems in Low-Dimensional Topology. Catalogue identity and source transcription; CC BY 4.0.
\url{https://huggingface.co/datasets/ulamai/UnsolvedMath}.
\item[S2] K3: A New Problem List in Low-Dimensional Topology, Problem 1.30. Expanded formulation of the supplied catalogue statement.
\url{https://math.berkeley.edu/sites/default/files/surv-295-ruberman-watermarked-author-pdf.pdf}.
\end{itemize}
\end{document}
